% =====================================================================
%  Resumo (NBR 6028) — em um relatorio, faz as vezes de sumario
%  executivo: quem le so esta pagina precisa entender o que foi feito,
%  o que foi encontrado e o que se recomenda.
%
%  As palavras-chave saem do latexkit.config.json.
% =====================================================================

\begin{resumo}
  Escreva aqui o resumo do relatorio: o objetivo do trabalho, o metodo
  empregado, os principais resultados obtidos e as recomendacoes
  decorrentes. Quem so ler esta pagina deve sair sabendo o essencial.

  \vspace{\onelineskip}
  \noindent
  \textbf{Palavras-chave}: \lkKeywords
\end{resumo}

% Abstract: obrigatorio em relatorios de pesquisa, dispensavel em
% relatorios internos. Apague o bloco abaixo se nao se aplicar.
\begin{resumo}[Abstract]
  \begin{otherlanguage*}{english}
    Write here the English version of the summary, keeping the same
    structure: purpose, method, results and recommendations.

    \vspace{\onelineskip}
    \noindent
    \textbf{Keywords}: \lkKeywordsEn
  \end{otherlanguage*}
\end{resumo}
